\chapter{绪论}
\label{ch:intro}

\section{研究背景}

蒸箱把饱和蒸汽送进封闭腔体，靠冷凝潜热把热量从外向内送进食品；烤箱用干热空气和管壁辐射把表面推到褐变温度，内部则靠导热慢慢跟上。两者分体布置时，厨房被两台机器占用，蒸制与烤制各自保留一套互不引用的定时曲线。蒸制时蒸汽量大、时间难校准，容易出现中心未熟或表面过湿；烤制时腔温波动、搁架与门缝造成边界不对称，内部冷点偏离几何中心\cite{feyissa2013}。漏电、干烧和超温是两类机器共有的安全问题，现有保护回路读取水位、门开关、剩余电流和加热器温度，与食品内部的进度估计分属两套回路。

家用蒸箱的常见容积在 20--30 L，饱和蒸汽由底部水箱或外置发生器供给，冷凝热流密度可以达到每平方米数千瓦，因此薄件表面很快升到接近饱和温度，内部仍按导热时间尺度慢慢跟上。对流烤箱的管壁辐射和热风把表面送到 \(150^\circ\mathrm{C}\) 以上，褐变和结壳发生在近表面几毫米，冷点仍由厚度平方决定。两台机器各自标定一套菜单：蒸箱按重量或按“鱼、蔬菜、馒头”给出分钟数，烤箱按腔温设定加定时。菜单一旦写进面板，这一次进门的初温、脂肪夹层和真实厚度都不再改写停机时刻。

厨房里的操作顺序把这件事写得更具体。一块从冰箱取出的肉糜，表面已经结了一层薄霜，中心仍接近冷藏温度；同一天里另一块室温面团，表面已经偏干。面板上的“蒸 20 分钟”或“烤 \(180^\circ\mathrm{C}\)、25 分钟”对这两块食品给出同一条时间轴。操作者看见的是蒸汽从门缝逸出、管壁发红、风机噪声升高，看不见的是冷点有没有越过 \(70^\circ\mathrm{C}\)、表面含水有没有掉到结壳门槛以下。课程任务把这一看不见的内部写成信息感知：观测通道能安装在门、搁架和管壁上，决策需要的却是厚度内部的热阻和进度。

把蒸和烤并入同一腔体，台面占用的问题当场消失。并入之后，操作者在面板上选择的是蒸、烤或组合，机器真正能改的仍只有四个量：加热器功率、蒸汽阀、排汽和风机。食品一旦进门，冷点温度、内部含水和反应进度都不在这四个旋钮上。现有商品路径——机械定时、腔温闭环、按重量调用曲线、穿刺探针、表面红外——闭合的是空气温度、申报质量或刺入点温度，止点写在第\ref{ch:lit}章。这一块食品的内部热阻，在这一次过程里从未进入后验。蒸汽把表面打湿，加热把表面推干，同一阀门对冷点和表面品质的符号可以相反：一体机把两个边界并进同一套旋钮之后，这个符号冲突变得无法回避。

课程任务要求以信息感知为基础设计智能蒸烤一体机。本文把感知写成：在可安装观测之下，决策泛函的剩余不确定度是否仍高于允许误判；若高于，执行器必须离开平衡去购买信息。腔温采样一旦进入饱和段，再加密同一通道也只是重复已经看见的边界量。操作者在面板上读到的“智能”，多数是把已经闭合的空气温度再画成一条更密的曲线。曲线再密，这一块食品的内部热阻仍然缺席。

家用过程的时间上限把这件事钉死。55 min 对应常见蒸烤时域，也对应数值核的积分终点。15 mm 的 Fourier 数在开环菜单 900 s 已经接近 1，冷点有机会跟上表面；50 mm 在同一上限内 Fourier 数只有 0.12，导热时间尺度还没有走完。于是同一台机器、同一套四个旋钮，薄件被菜单推过设计线，厚件停在设计线之下。差别首先来自厚度，其次才来自策略名称。第\ref{ch:result}章把这一几何事实写成交叉表，第\ref{ch:wide}章把 50 mm 时域放到 70 min 之后才看见同一初值走到 \(\alpha_m=0.85\)。

课程报告因此收成一台可以复核的机器。同一积分器、同一套四个执行量、同一条设计线 \(\alpha_m=0.85\)，开环、腔温比例、无探针反馈和对偶四条策略共用这条线。学号 20241334、姓名庄颂、班级信息2402 写在封面与页眉，正文随后给出的交叉表、网格加密、因子消融和种子 7 轨迹族，比较的都是这四条策略在十六条控制体上的停机偏差。交互页把同一条状态流读回三维腔体，静态地址为 \url{https://2711944586.github.io/xxgz/}。

\section{问题陈述}

记食品与腔体的真实状态为 \(\vect{x}(t)\)，可安装传感器给出
\[
\vect{y}(t)=h(\vect{x}(t))+\vect{v}(t).
\]
停机、表面品质和安全构成决策泛函 \(\phi(\vect{x})\)。沿 Lindley 的实验信息定义\cite{lindley1956}，信息赤字为
\begin{equation}
\label{eq:deficit}
\Delta\mathcal{I}(t)=H\big(\phi(\vect{x}(t))\mid \vect{y}_{0:t}\big)-H_{\mathrm{tol}}.
\end{equation}
\(H_{\mathrm{tol}}\) 由动作切换的滞回带决定。\(\Delta\mathcal{I}>0\) 时若仍按开环时间表改边界，冷点进度和表面干耗会分开漂移：二者对同一蒸汽指令的符号不必相同。

因此机器要完成的链是
\begin{equation}
\label{eq:chain}
\begin{aligned}
&\text{边界激励}\to\text{多尺度状态}\to\text{后验}\\
&\to\text{信息价值与品质价值}\to\text{安全投影}\to\text{新的边界}.
\end{aligned}
\end{equation}
链上缺一环，后面的环就没有输入。没有后验，信息价值和品质价值无法分开；没有安全投影，品质模型的错误会直接变成执行器指令。

图\ref{fig:cavity}给出本文的工作对象。加热器在顶部，蒸汽从一侧进入，风机带动气流，排汽带走湿空气。档位只限制哪几个量允许离开零：蒸档把加热请求夹在 0.25 以下，烤档把蒸汽写成零并把排汽抬到至少 0.45，组合档放开四个量。水位、门、剩余电流和加热器温度是硬测量：其中任何一条越过阈值，加热和蒸汽同时写成零。进度、热阻和估计冷点留在品质回路里，硬切断的自变量只有这四条测量。

家用过程还有一条时间上限。本文把时域写成 55 min，对应家用蒸烤的常见上限，也对应数值核的积分终点。厚件在这条上限内走不完导热时间尺度，这一事实会在第\ref{ch:result}章和第\ref{ch:wide}章反复出现：探针可以把热阻方差降下来，扩散时间仍把 50 mm 钉在设计线之下。

\begin{figure}[htbp]
\centering
\includegraphics[width=0.88\textwidth]{cavity_schematic.png}
\caption{同一腔体上的四个执行量。食品沿厚度离散为十六条控制体，冷点取最低温度，不必落在几何中心。}
\label{fig:cavity}
\end{figure}
\FloatBarrier

\section{研究目标}

本文要同时满足五条性质。只成立其中几条时，系统会退回到已经有的装置：只会定时，只会跟腔温，或只会在超温时断电。

\begin{enumerate}
\item \textbf{可重构。} 冷点温度、表面含水、反应进度和内部热阻能由有限观测与短时边界激励更新。稳态观测不含内部热阻，可重构依赖离开平衡的激励。
\item \textbf{可分离。} 同一次阀门动作要么为了降低后验熵，要么为了推动品质泛函。日志里要能查出这一时刻是哪一个函数发出的指令。
\item \textbf{可迁移。} 宏观传热、孔隙水分和反应进度耦合成低维决策状态。食材类别只改变先验和终端目标，不改变状态的分量定义。
\item \textbf{可更新。} 单次过程只更新事先声明的传递参数。结构方程的父节点不随这次数据改写。
\item \textbf{可拦截。} 品质模型的输出在进入执行器之前，被只依赖硬测量的不等式投影。硬保护不经过该模型。
\end{enumerate}

这五条是设计判据。已经运行的一维核把可分离落实为 150--190 s 的固定窗：窗内蒸汽与风机按时间表离开平衡，窗外加热请求由冷点估计、方差和厚度决定。热阻更新使用由真值温差构成的参照，估计器同时读取可安装通道和评分参照；把参照换成可安装通道，是把可重构写严的第一步，见第\ref{ch:conc}章。安全投影落实到加热器温度这一路；满功率平衡点 \(196.5^\circ\mathrm{C}\) 低于 \(230^\circ\mathrm{C}\) 阈值，削减在战役里保持待命，规则只读管温。完整读法写在第\ref{ch:conc}章。

五条性质一起成立，机器才从“会定时的腔温仪表”变成“这一次过程里看见内部热阻的蒸烤一体机”。已经运行的一维核把可拦截落实为加热器温度这一路，把可迁移落实为九维决策状态的分量定义。可重构和可更新写得更严的步骤，放在第\ref{ch:conc}章：热阻更新目前仍读取由真值温差构成的参照，战役结束于日志，后验写回下一次先验留在规格里。

\section{主要工作}

\begin{enumerate}
\item 把蒸烤一体机写成四个边界上的信息问题。文献中的体积平均、品质动力学、逆热传导、对偶控制和屏障函数被标明继承层与规格层。
\item 给出从多孔介质连续方程到十六条控制体、再到两节点工作方程的收缩，并定义九维决策状态。
\item 证明平衡观测对内部热阻灵敏度为零，两段相变变号加对流的探针提供第二个灵敏度方向，并给出探针窗口内冷点增量的上界。
\item 把同一次蒸汽分解为冷点潜热、表面蒸发和结壳热阻三条路径，安全投影的可行域与品质后验分离。
\item 在同一数值核上比较开环、腔温比例、无探针反馈和对偶四种策略，并给出蒸、烤、组合三档的交叉结果。
\item 实现读取同一状态流的三维腔体演示，边界、色带、路径和投影同屏可见。
\end{enumerate}

\section{论文结构}
\enlargethispage{3\baselineskip}

第\ref{ch:lit}章评述文献与产品路径。第\ref{ch:arch}章给出闭环与三档允许集合。第\ref{ch:physics}章建立多尺度模型。第\ref{ch:est}章讨论观测与估计。第\ref{ch:ati}章给出探针与四种策略。第\ref{ch:safe}章处理因果路径与安全投影。第\ref{ch:result}章报告主对照与档位研究。第\ref{ch:grid}章做网格收敛。第\ref{ch:ablate}章做因子消融与边界扰动。第\ref{ch:wide}章给出宽格子、集合散布与 50 mm 时域延长。第\ref{ch:section}章处理二维截面、通道关联与 Fourier--Biot 尺度。第\ref{ch:atlas}章把交叉表读回种子 7 的轨迹族。第\ref{ch:demo}章说明三维演示与一键启动。第\ref{ch:conc}章总结。符号、离散格式、种子 7 轨迹和尺度估计放在附录。

阅读顺序可以按问题走：若先问“机器能改什么”，从第\ref{ch:arch}章的四个执行量和三档允许集合开始；若先问“内部热阻为什么看不见”，从第\ref{ch:est}章的平衡命题和半无限体放大表开始；若先问“数字从哪来”，从第\ref{ch:result}章的 S1 与档位研究开始，再回到第\ref{ch:ati}章的时间表。演示章放在结论之前，用来把交叉表读回可见零件。

课程报告因此同时给出三件可以独立复核的工件。数值核沿厚度用十六条控制体，时间步 \(2\,\mathrm{s}\)，时域 55 min，主对照 19440 行、档位研究 11664 行，种子 7 的轨迹写进附录；XeLaTeX 文稿按问题组织，封面与页眉标明学号 20241334、姓名庄颂、班级信息2402；浏览器演示读取同一状态流，三维腔体里六根加热管、十叶风机、蒸汽粒子和十六块薄片与日志对齐到写入精度。三件工件共用种子、厚度和策略名称，求解器只在 Python 核里。公式编号以 PDF 为准。
